% Working fragment; finalize bibliography keys and full-study eligibility before integration.
\subsection{Evidence needed for learning-assisted congestion control}
Learning-assisted congestion control covers different control actions. A policy
may select an existing congestion-control algorithm, adjust the congestion
window directly, or modify a parameter within a BBR implementation. These
actions have different feedback and deployment requirements, so performance
observations should be compared within the evaluated control task.

For example, the inspected WiFi study (DOI 10.3390/fi16090330) uses offline
learning to select among congestion-control algorithms. Its stated reward
combines throughput, latency and loss. This provides evidence about the
reported selection setup, but does not by itself establish a safe online
BBRv3 parameter controller or an inference-time budget. TCP-Drinc
(DOI 10.1109/ACCESS.2019.2892046), in contrast, learns congestion-window
actions in ns-3 using a goodput--delay utility difference. Its fairness table
reports a mean Jain index of 0.8058 with a 95\% confidence interval of
[0.7233, 0.8882]. These results remain specific to the simulation conditions;
the reported baseline intervals overlap, and a significant improvement in
fairness is not established by the table alone.

Training accuracy, accumulated reward and a narrow reported confidence
interval answer different questions. Deployment assessment requires measured
inference latency and resource use, while transfer assessment requires tests
outside the training conditions. Fairness assessment should report the
objective, per-flow rates and the evaluation metric together. A reward need
not contain Jain's index explicitly to express a fairness-related utility;
conversely, training convergence alone does not prove fairness under a new
queue policy, flow population or implementation.
